\documentclass[11pt,a4paper]{article}
\usepackage[T1]{fontenc}
\usepackage{lmodern}
\usepackage[utf8]{inputenc}
\usepackage[polish]{babel}
\usepackage{amsmath,amssymb,amsthm,hyperref}
\usepackage[margin=2.5cm]{geometry}
\newtheorem{theorem}{Twierdzenie}
\newtheorem{lemma}[theorem]{Lemat}
\theoremstyle{definition}
\newtheorem{definition}{Definicja}[section]
\numberwithin{equation}{section}
\newcommand{\norm}[1]{\left\lVert #1\right\rVert}
\title{Rozwiązania listy 12 z przedmiotu Programy użytkowe}
\author{Przykładowe opracowanie}
\date{\today}
\begin{document}
\maketitle
\noindent\textbf{Źródło zadań:} plik \texttt{old/lista12.tex} w tym
repozytorium. Jest to archiwalna lista; aktualnej listy student nie
dostarczył. Rozwiązanie pokazuje komplet techniczny trzech poleceń.

\section{Środowiska, liczniki i przykład użycia}
\noindent\textbf{Treść zadania 1.} Zdefiniuj środowiska Definicja,
Twierdzenie i Lemat. Twierdzenia oraz lematy mają wspólny licznik.
Definicje i równania numeruj w obrębie sekcji.

\medskip
\noindent\textbf{Rozwiązanie.} W preambule \texttt{lemma} dzieli licznik
z \texttt{theorem}, zaś \texttt{definition} ma licznik resetowany
na początku sekcji. Polecenie \verb|\numberwithin{equation}{section}|
ustawia numer równania w postaci numer sekcji.numer równania.
\begin{definition}
Normą euklidesową wektora $x=(x_1,\ldots,x_n)$ jest
$\norm{x}_2=\sqrt{\sum_{i=1}^{n}x_i^2}$.
\end{definition}
\begin{definition}
Iloczynem skalarnym wektorów $x,y\in\mathbb R^n$ jest
$\langle x,y\rangle=\sum_{i=1}^{n}x_i y_i$.
\end{definition}
\begin{lemma}
Dla każdego $x\in\mathbb R^n$ mamy $\norm{x}_2\geq 0$.
\end{lemma}
\begin{proof}
Każdy składnik $x_i^2$ jest nieujemny, a pierwiastek
z liczby nieujemnej jest nieujemny.
\end{proof}
\begin{theorem}
Dla $x,y\in\mathbb R^n$ zachodzi nierówność trójkąta
$\norm{x+y}_2\leq\norm{x}_2+\norm{y}_2$.
\end{theorem}
\begin{proof}
Z nierówności Cauchy'ego--Schwarza otrzymujemy
\begin{equation}
\norm{x+y}_2^2
=\norm{x}_2^2+2\langle x,y\rangle+\norm{y}_2^2
\leq(\norm{x}_2+\norm{y}_2)^2.
\end{equation}
Pierwiastkowanie obu nieujemnych stron kończy dowód.
\end{proof}
Numeracja lematów i twierdzeń jest wspólna, a definicje i równania
mają numery przypisane do sekcji.

\section{Własna komenda}
\noindent\textbf{Treść zadania 2.} Zdefiniuj własną komendę z argumentem,
która upraszcza zapis wyrażenia.

\medskip
\noindent\textbf{Rozwiązanie.} W preambule zdefiniowano
\verb|\norm{#1}|. Użycie \verb|\norm{x-y}| daje $\norm{x-y}$.
Komenda opakowuje argument w skalowalne znaki normy, więc
unikamy powtarzania dwóch poleceń: \verb|\left\lVert|
oraz \verb|\right\rVert|. Przykład:
\begin{equation}
\norm{x-y}_2^2=\sum_{i=1}^{n}(x_i-y_i)^2.
\end{equation}

\section{Bibliografia}
\noindent\textbf{Treść zadania 3.} Utwórz plik z przykładową
bibliografią i dołącz go do dokumentu za pomocą BibTeX-a
lub \texttt{biblatex}.

\medskip
\noindent\textbf{Rozwiązanie.} Plik \texttt{references.bib} zawiera wpis
książkowy. Cytowanie w tekście: \cite{lamport1994}.
Kompilacja: \texttt{pdflatex}, \texttt{bibtex},
\texttt{pdflatex}, \texttt{pdflatex}.
\bibliographystyle{plain}
\bibliography{references}
\end{document}
